\clearpage\pagestyle{otMain}\setlength{\otBlocEspacement}{2.5mm}\fontsize{11}{15}\selectfont\corrigesChapitre{8}{Pourcentages et grandeurs quotients}
\begin{blocLivre}{cor-8-1}\begin{corrige}{1}{Une part en pourcentage}\demarche{Multiplier la quantité par la fraction correspondante.}
$18\,\%=\frac{18}{100}$. La part cherchée vaut $\frac{18}{100}\times250=\res{45}$.
\end{corrige}\end{blocLivre}
\begin{blocLivre}{cor-8-2}\begin{corrige}{2}{Retrouver le pourcentage}\demarche{Diviser la partie par le total.}
Le rapport partie sur total est $\frac{12}{32}=0{,}375$. Il correspond à $\res{37{,}5\,\%}$. Le dénominateur est l’effectif total, pas celui des autres élèves.
\end{corrige}\end{blocLivre}
\begin{blocLivre}{cor-8-3}\begin{corrige}{3}{Des coefficients}\demarche{Ajouter ou retrancher la part en centièmes à 1.}
Hausse : $1+\frac{t}{100}$ ; baisse : $1-\frac{t}{100}$. On obtient $\res{1{,}15}$, $\res{0{,}8}$, $\res{1{,}03}$ et $\res{0{,}925}$. Un coefficient inférieur à 1 traduit une baisse.
\end{corrige}\end{blocLivre}
\begin{blocLivre}{cor-8-4}\begin{corrige}{4}{Une vitesse moyenne}\demarche{Exprimer la durée dans l’unité du quotient.}
$2$ h $30$ min $=2{,}5$ h. La vitesse moyenne est $45:2{,}5=\res{18\ \mathrm{km/h}}$.
\end{corrige}\end{blocLivre}
\begin{blocLivre}{cor-8-5}\begin{corrige}{5}{Une durée}\demarche{Utiliser une heure de 60 minutes.}
$18$ min $=\frac{18}{60}=0{,}3$ h, donc la durée vaut $\res{2{,}3\ \mathrm h}$. En minutes : $2\times60+18=\res{138}$ min.
\end{corrige}\end{blocLivre}
\begin{blocLivre}{cor-8-6}\begin{corrige}{6}{Une remise}\demarche{Multiplier directement par le coefficient.}
Le coefficient vaut $1-0{,}12=0{,}88$. Le nouveau prix est $80\times0{,}88=\res{70{,}40}$ euros. La remise vaut 9,60 euros.
\end{corrige}\end{blocLivre}
\begin{blocLivre}{cor-8-7}\begin{corrige}{7}{Une hausse de loyer}\demarche{Associer la hausse au coefficient supérieur à 1.}
Le coefficient vaut $1+0{,}025=1{,}025$. Le nouveau loyer est $640\times1{,}025=\res{656}$ euros. La hausse vaut $656-640=\res{16}$ euros, soit $\frac{2{,}5}{100}\times640$.
\end{corrige}\end{blocLivre}
\begin{blocLivre}{cor-8-8}\begin{corrige}{8}{Une conversion de vitesse}\demarche{Convertir distance et durée.}
72 km correspondent à $72\,000$ m en 3 600 s. Donc $72\,000:3\,600=\res{20\ \mathrm{m/s}}$. On retrouve $72:3{,}6=20$.
\end{corrige}\end{blocLivre}
\begin{blocLivre}{cor-8-9}\begin{corrige}{9}{Changer les deux unités}\demarche{Convertir séparément le volume et le temps de référence.}
On convertit d’abord le volume : $0{,}012$ m$^3$/s $=12$ L/s. En une minute, il s’écoule $12\times60=\res{720\ \mathrm{L/min}}$. Une minute contient davantage de volume qu’une seconde.
\end{corrige}\end{blocLivre}
\begin{blocLivre}{cor-8-10}\begin{corrige}{10}{Une masse volumique}\demarche{Calculer le quotient puis réutiliser le coefficient.}
$\rho=\frac{108}{40}=\res{2{,}7\ \mathrm{g/cm}^3}$ : chaque cm$^3$ pèse 2,7 g. À masse volumique identique, $m=2{,}7\times65=\res{175{,}5}$ g.
\end{corrige}\end{blocLivre}
\begin{blocLivre}{cor-8-11}\begin{corrige}{11}{Deux remises successives}\demarche{Multiplier les coefficients successifs.}
Les coefficients se multiplient : $100\times0{,}8\times0{,}9=\res{72}$ euros. Le coefficient global $0{,}72$ correspond à une baisse de $\res{28\,\%}$, et non de 30\,\% : la seconde baisse porte sur 80 euros.
\end{corrige}\end{blocLivre}
\begin{blocLivre}{cor-8-12}\begin{corrige}{12}{Un pourcentage inverse}\demarche{Diviser par le coefficient multiplicateur.}
Le prix payé est le prix initial $P$ multiplié par 0,75 : $0{,}75P=54$, d’où $P=54:0{,}75=\res{72}$ euros. Vérification : 25\,\% de 72 vaut 18, et $72-18=54$.
\end{corrige}\end{blocLivre}
\begin{blocLivre}{cor-8-13}\begin{corrige}{13}{Deux groupes réunis}\demarche{Revenir aux effectifs avant de réunir les groupes.}
Il y a $20\times0{,}4=8$ filles et $30\times0{,}6=18$ femmes, soit 26 personnes sur 50. Le pourcentage est $\frac{26}{50}=\res{52\,\%}$. La moyenne simple de 40\,\% et 60\,\% ignorerait les effectifs différents.
\end{corrige}\end{blocLivre}
\begin{blocLivre}{cor-8-14}\begin{corrige}{14}{Aller et retour}\demarche{Diviser la distance totale par la durée totale.}
L’aller dure $30:15=2$ h et le retour $30:30=1$ h. Pour 60 km en 3 h, la vitesse moyenne vaut $\res{20\ \mathrm{km/h}}$. La moyenne simple 22,5 ne tient pas compte des durées différentes.
\end{corrige}\end{blocLivre}
\begin{blocLivre}{cor-8-15}\begin{corrige}{15}{Une comparaison de débits}\demarche{Comparer dans une même unité puis calculer les durées.}
B débite $0{,}018\times1\,000\times60=\res{1\,080}$ L/min, donc plus que A. A met $5\,400:900=6$ min ; B met $5\,400:1\,080=5$ min. Le gain est de $\res1$ min.
\end{corrige}\end{blocLivre}
